% FROZEN — do not keep editing. Living source is harsh/resume/HarshDesai_Resume.tex
% This is the Terafab/plasma recipe that was compiled (2026-07-28).
% Content draft that fed this file: Harsh_Desai_Terafab_Resume.tex (same folder).
% PDFs from this file: HarshDesai_Terafab_PlasmaSystems_Resume.pdf and HarshDesai_Resume_July2026_Plasma.pdf (duplicate).
\documentclass[10pt,a4paper]{article}

\usepackage[a4paper,top=0.31in,bottom=0.28in,left=0.39in,right=0.39in]{geometry}
\usepackage{fontspec}
\setmainfont{Tinos}
\usepackage{xcolor}
\usepackage{enumitem}
\usepackage{tabularx}
\usepackage{array}
\usepackage{hyperref}
\usepackage{titlesec}
\usepackage{microtype}
\usepackage{ragged2e}

\definecolor{nameblue}{HTML}{214F78}
\definecolor{linkblue}{HTML}{0000CC}

\hypersetup{
  colorlinks=true,
  urlcolor=linkblue,
  linkcolor=linkblue,
  pdfauthor={Harsh Desai},
  pdftitle={Harsh Desai Resume}
}

\pagestyle{empty}
\setlength{\parindent}{0pt}
\setlength{\tabcolsep}{0pt}
\setlength{\emergencystretch}{1em}
\fontsize{8.18pt}{8.78pt}\selectfont

\titleformat{\section}
  {\fontsize{10.2pt}{10.8pt}\selectfont\bfseries\uppercase}
  {}{0pt}{}
  [\vspace{-2.2pt}\titlerule]
\titlespacing*{\section}{0pt}{6.0pt}{1.6pt}

\setlist[itemize]{
  leftmargin=13.5pt,
  label=\raisebox{0.18ex}{\tiny$\bullet$},
  itemsep=0.15pt,
  parsep=0pt,
  topsep=0.25pt,
  partopsep=0pt
}

\newcommand{\resumeSubheading}[4]{%
  \begin{tabular*}{\textwidth}{@{\extracolsep{\fill}}l r}
    \textbf{#1} & \textcolor{nameblue}{\textbf{#2}} \\
    #3 & #4
  \end{tabular*}\vspace{-1.6pt}
}

\newcommand{\resumeProjectHeading}[2]{%
  \begin{tabular*}{\textwidth}{@{\extracolsep{\fill}}l r}
    \textbf{#1} & #2
  \end{tabular*}\vspace{-1.6pt}
}

\newcommand{\resumeItemsStart}{\begin{itemize}}
\newcommand{\resumeItemsEnd}{\end{itemize}\vspace{-1.1pt}}
\newcommand{\resumeItem}[1]{\item #1}

\begin{document}

% Header
\begin{tabular*}{\textwidth}{@{\extracolsep{\fill}}l r}
  {\fontsize{18pt}{18pt}\selectfont\bfseries\textcolor{nameblue}{Harsh Desai}} &
  \fontsize{8.05pt}{8.5pt}\selectfont
  Ann Arbor, MI, USA $\mid$
  \href{mailto:harshdes@umich.edu}{harshdes@umich.edu} $\mid$
  +1-734-548-1080 $\mid$
  \href{https://www.linkedin.com/in/harshddes/}{LinkedIn} $\mid$
  \href{https://harshddes.github.io/}{harshddes.github.io}
\end{tabular*}
\vspace{-3.5pt}

{\itshape Experimental systems engineer in low-pressure plasma, vacuum/HV test equipment, automation, detector calibration, and ADC/FPGA readout; targeting semiconductor process, equipment, and diagnostics roles.}
\vspace{0.4pt}

\section{Education}
\resumeSubheading
  {University of Michigan}{Ann Arbor, MI, USA}
  {Master of Engineering, Space Systems Engineering $\mid$ GPA: 3.79/4.0}{Aug 2024--Dec 2025}
\resumeSubheading
  {Vellore Institute of Technology (VIT)}{Vellore, India}
  {Bachelor of Technology, Mechanical Engineering $\mid$ GPA: 7.78/10}{Jul 2019--Jul 2023}

\section{Technical Skills}
\textbf{Plasma, Vacuum \& Process Equipment:} Low-pressure hollow-cathode plasma operation, 1.3 kV HV interface validation, vacuum chamber operation (CTI-Cryogenics Hi-Vac; roughing pump operated at 1 mTorr), RPA, emissive-probe and Langmuir-probe diagnostics, CEM/ESA calibration\\[-0.2pt]
\textbf{Automation, Controls \& DAQ:} Python, LabVIEW (CLAD prep), MATLAB (learning), GPIB/IEEE-488, USB-serial, serial data logging, NI VISA (PyVISA, PyMeasure), Keithley DAQs and SMUs, TDK Lambda PSUs, High Voltage Switching Operations\\[-0.2pt]
\textbf{FPGA, Digital Systems \& Modeling:} Verilog (learning), VHDL (learning), SIMION, SRIM, Amptek DPPMCA, FPGA Design Workflow, AMD Vivado, MATLAB HDL Coder \& DSP Toolbox, SPENVIS, SWMF \& Magnetosphere-Ionosphere Models (MAGE, SWMF, HYPERS), Google Data Studio\\[-0.2pt]
\textbf{Mechanical \& Manufacturing Tools:} ANSYS (GUI/TUI): Fluent, Mechanical; PyANSYS, OptiSlang, OpenFOAM, SolidWorks, Fusion 360, AutoCAD, Autodesk Inventor, SpaceClaim, MIDO

\section{Research \& Engineering Experience}

\resumeSubheading
  {Climate and Space Sciences and Engineering, University of Michigan}{Ann Arbor, MI, USA}
  {Research Assistant, LVACCS Testing Rig Characterization}{Feb 2026--Present}
\resumeItemsStart
  \resumeItem{Validated a 1300 V hollow-cathode plasma-source test workflow by combining HV discharge-box FMEA, Python/PyVISA ignition control, PSU logging, and DAQ synchronization; cut manual steps 15$\rightarrow$5 and achieved 98.6\% synchronized data}
  \resumeItem{Supported electrostatic regolith-lofting experiments using Retarding Potential Analyzer (RPA), emissive-probe, and Langmuir-probe diagnostics to compare plasma conditions across test states}
  \resumeItem{Characterized remote ignition and post-run processing workflows for Spacecraft Charging Device test readiness, improving experimental repeatability and saving $\sim$15 minutes of processing per run}
  \resumeItem{Built practical experience with low-pressure discharge operation, ignition stability, chamber diagnostics, and parameterized testing, creating a technical bridge toward plasma etch and thin-film deposition equipment}
\resumeItemsEnd

\resumeSubheading
  {Space Physics Research Lab, University of Michigan}{Ann Arbor, MI, USA}
  {Summer Intern, TestBedz Spacecraft Qualification Platform}{May 2025--Jul 2025}
\resumeItemsStart
  \resumeItem{Built a full-stack requirement-flowdown platform for TVAC, vibration, and EMI test submissions, matching spacecraft hardware profiles to facility limits across 6 test-type workflows}
  \resumeItem{Encoded hardware metadata, facility operating envelopes, and acceptance criteria into machine-readable compatibility checks and traceable test plans for repeatable equipment-qualification workflows}
\resumeItemsEnd

\resumeSubheading
  {Solar and Heliospheric Research Group, University of Michigan}{Ann Arbor, MI, USA}
  {Student Research Assistant, SSD Readout Chain}{Jan 2025--Dec 2025}
\resumeItemsStart
  \resumeItem{Designed an SSD readout roadmap around science-driven energy-resolution targets, integrating a Zmod ADC 1410 path with a Zynq-7000 architecture to evaluate a 1 MSPS$\rightarrow$125 MSPS digitization upgrade}
  \resumeItem{Verified MATLAB HDL Coder golden models against Vivado co-simulations and diagnosed 125 MHz post-synthesis timing limits, including Worst Negative Slack constraints in the Zynq SoC signal path}
\resumeItemsEnd

% \resumeSubheading
%   {Climate and Space Sciences and Engineering, University of Michigan}{Ann Arbor, MI, USA}
%   {Research Assistant, Uranian Tropospheric Cloud Resolving Model}{Sep 2025--Dec 2025}
% \resumeItemsStart
%   \resumeItem{Ran Uranian cloud-resolving model sensitivity studies on Great Lakes HPC, linking methane profile, microphysics, latent heating, gravity, and sedimentation effects to vertical-velocity and condensation diagnostics}
% \resumeItemsEnd

\resumeSubheading
  {L3Harris and University of Michigan}{Ann Arbor, MI, USA}
  {Communications Lead, Uranian Orbiter and Probe Mission Study}{Sep 2024--May 2025}
\resumeItemsStart
  \resumeItem{Led Uranian mission communications architecture trades and DSN/fault-tolerance requirement flowdown, managing 15 traceable requirements and supporting a 16\% (\$450M) cost reduction through an orbiter-only configuration trade}
\resumeItemsEnd

\section{Selected Instrumentation \& Process-Relevant Projects}
\resumeProjectHeading{Space Instrumentation Calibration \& Ion-Optics Series}{Jul 2025--Dec 2025}
{\itshape Graduate Course: SPACE 571 $\mid$ Supervisor: Prof. Stefano Livi}\vspace{-1.3pt}
\resumeItemsStart
  \resumeItem{Calibrated Channel Electron Multiplier response by mapping operating bias, pulse-height distributions, and amplifier gain behavior, connecting detector settings to usable signal quality and count-rate interpretation}
  \resumeItem{Calibrated the signal chain from charge injection through Charge-Shaper Amplifier, pulse shaping, and MCA acquisition, extracting gain, centroid, FWHM, and bias-dependent count response}
  \resumeItem{Modeled Electrostatic Analyzer and ion-optics performance using lab energy-angle scans, SIMION, and SRIM to quantify transmission, angular acceptance, energy resolution, FWHM, and filter response}
\resumeItemsEnd

% \resumeProjectHeading{Spatial Ion-Electron Temperature Correlation}{Jan 2025--Apr 2025}
% {\itshape Graduate Course: SPACE 477 $\mid$ Supervisor: Prof. Xianzhe Jia}\vspace{-1.3pt}
% \resumeItemsStart
%   \resumeItem{Analyzed ion-electron temperature structure in SWMF/HYPERS simulations, tracing magnetosheath ion distributions and reconnection-driven electron heating during simulated shock events}
% \resumeItemsEnd


\section{Projects \& Leadership}
\resumeProjectHeading{CANSAT 2022 (NASA \& AAS) $\mid$ Ranked 7th worldwide out of 42 teams}{Aug 2021--Jul 2022}
\resumeItemsStart
  \resumeItem{Built a 10 m tether-deployment, payload release, and camera-stabilization system using a DC motor, worm gear, torsion-spring latch, and dual-servo gimbal for controlled descent and fixed 45$^{\circ}$ imaging}
\resumeItemsEnd

\resumeProjectHeading{Spaceport America Cup / IREC $\mid$ Placed 23rd globally and 5th in Asia-Pacific}{Apr 2020--Jul 2021}
\resumeItemsStart
  \resumeItem{Launched a Mach 0.9 solid-motor rocket to 10,000 ft and developed an SRAD trajectory simulator using CFD-backed aerodynamic data for flight prediction}
\resumeItemsEnd

\resumeProjectHeading{CANSAT 2021 (NASA \& AAS) $\mid$ Placed 13th globally and 7th in Asia-Pacific}{Aug 2020--Jul 2021}
\resumeItemsStart
  \resumeItem{Modeled maple-seed mono-wing payload descent using Blade Element Theory and transient CFD, with MATLAB telemetry software for real-time dual-payload deployment monitoring}
\resumeItemsEnd

\end{document}
